\scanpage{35}

et $H^0$ sont exacts à gauche, on en conclut un homomorphisme de diagrammes exacts d'ensembles
% No arrow strokes are visible between these six objects.
\diagram{
\begin{tikzcd}[ampersand replacement=\&,column sep=2.0em,row sep=1.7em]
H^0(X,F) \& H^{0}(X,G) \& H^0(X,H) \\
H^0(Y,F_Y) \& H^0(Y,G_Y) \& H^0(Y,H_Y),
\end{tikzcd}
}
où d'après ce qui précède les flèches verticales sont injectives\inserted{(ennparticulier la dernière}, et où par hypothèse la flèche verticale médiane est bijective. (NB les indices $Y$ désignantt les restrictions à $Y$)? On en conclut formellement que la première flèche verticale est également bijective.
% The inserted parenthesis runs into the right edge; no closing parenthesis is visible.

Nous appliquerons cette remarque au cas \deleted{où}
\[
G=h\ h\ (F)\ x\ (\quad(F))\ ,
\]
% Spaces in this and the following formulas reproduce unfilled places in the typescript.
et à l'homomorphisme évident $F\qquad G$, déduit des homomorphismes naturels dans les deux facteurs. Cet homomorhisme est injectif, comme il résulte immédiatement (par considération des fibres) de l'hypothèse $X=Y\quad(\overline{X})$. Il reste à prouver l'assertion de bijectivité pour $G$, ou ce qui revient au même, pour les deux facteurs, donc pour des faisceaux de la forme $h(\unsure{H|\mathbb{H}})$ (resp. $(\overline{F})$),où $H$ (resp. $\overline{F}$) est un faisceau sur $Y$ (resp. sur $\overline{X}$). or l'assertion est triviale pour un faisceau de la forme $h(H)$. Pour le cas d'un faisceau de la forme $(\overline{F})$, compte tenu du fait que l'homomorphisme de c\deleted{y}hangement de base $H(\quad(\overline{F}))\qquad{}_{Y}(\overline{h}(\overline{F}))$ est bijectif, on voit que l'homomorphisme $H^0(X,\quad(\overline{F}))\qquad H^0(Y,h\quad(\overline{F}))$ s'identifie à l'homomorphisme $H^0(\overline{X},\overline{F})\qquad H^0(\overline{Y},\overline{h}\ (\overline{F}))$, qui est bijectif par hypothèse. Cela achève la démonstration de (i).

(ii) La conclu\deleted{\unsure{t|b}}sion d\deleted{'}e\deleted{inject} bijectivité pour $H^0(X,F)\qquad H^0(Y,h\ (F))$ \margin{\handwritten{$\unsure{\backslash|\checkmark}$}}pour tout faisceau d'ensembles $F$ sur $X$, déjà obtenue dans (i), assure \deleted{par un argument immédiat} que pour tout faisceau en groupes $F$ sur $X$ (pas nécessairement ind-finis), l\deleted{'}e\deleted{homomorphisme du} foncteur restriction
% archive sheet number: 34
